% U3 WS1 -- Type B: IMFs, solids, phase behavior. Blocks 1-2.
\documentclass{apchem-worksheet}

\apcunit{3}
\apcunittitle{Properties of Substances and Mixtures}
\apcdoctitle{Worksheet 1 \textbullet\ Forces, Solids, Phases}
\apcsubtitle{CED 3.1--3.3 \textbullet\ every ranking needs its reason}

\begin{document}
\wsheader[Name the force AND the structural feature that produces it. ``It
has stronger IMFs'' restates the question --- say \emph{which} force and
\emph{why} it is present.]

\problem[3.1]{List every intermolecular force present in each substance:}
\begin{parts}
  \item \ce{CO2}: \answerblank[6cm]{LDF only (nonpolar)}
  \item \ce{CH3F}: \answerblank[6cm]{LDF + dipole--dipole}
  \item \ce{NH3}: \answerblank[7.1cm]{LDF + dipole--dipole + H-bonding}
  \item \ce{KBr} dissolved in water: \answerblank[6cm]{ion--dipole}
\end{parts}

\problem[3.1]{\ce{HF} boils at \SI{19.5}{\celsius} while \ce{HCl} boils at
\SI{-85}{\celsius}, even though \ce{HCl} has the greater molar mass.
Explain. \answerlines[3]{\ce{HF} hydrogen bonds --- H attached directly to
F, one of only three elements that permit it. That attraction far exceeds
the dipole--dipole and dispersion forces in \ce{HCl}, more than
compensating for \ce{HCl}'s larger electron count.}}

\problem[3.1]{Zumdahl reports these boiling points: \ce{BI3}
\SI{209.5}{\celsius} (LDF only), \ce{PCl3} \SI{76.1}{\celsius}
(LDF + dipole--dipole), \ce{NH3} \SI{-33.0}{\celsius} (LDF + H-bonding).}
\begin{parts}
  \item This ordering appears to contradict ``hydrogen bonding is
        strongest.'' Resolve the apparent contradiction.
        \answerlines[3]{Force type alone doesn't set the ranking ---
        magnitude does. \ce{BI3} has an enormous, highly polarizable
        electron cloud, so its dispersion forces are larger in absolute
        terms than \ce{NH3}'s hydrogen bonds. \ce{NH3} is also very small
        with few electrons.}
  \item State the rule this illustrates in one sentence.
        \answerlines[2]{Compare molar mass/polarizability first when
        molecular sizes differ greatly; only compare force types among
        molecules of similar size.}
\end{parts}

\problem[3.2]{Classify each solid and predict whether it conducts
electricity as a solid:}
\begin{parts}
  \item \ce{SiC}, mp \SI{2730}{\celsius}, extremely hard:
        \answerblank[5.5cm]{covalent network; no}
  \item \ce{Ag}, mp \SI{962}{\celsius}, malleable:
        \answerblank[5.5cm]{metallic; yes}
  \item \ce{K2S}, mp \SI{840}{\celsius}, brittle:
        \answerblank[5.5cm]{ionic; no (yes when molten)}
  \item \ce{I2}, mp \SI{114}{\celsius}, soft:
        \answerblank[5.5cm]{molecular; no}
\end{parts}

\problem[3.3]{A student claims that melting ice breaks the O--H bonds in
water. Evaluate this claim and give the correct account.
\answerlines[3]{False. Melting overcomes \emph{intermolecular} hydrogen
bonds between molecules; the covalent O--H bonds within each molecule are
untouched. Breaking O--H bonds would decompose water into different
substances and requires far more energy.}}

\problem[3.3]{Rank by vapor pressure at \SI{25}{\celsius}, highest first,
with a one-line reason: \ce{H2O}, \ce{CH3OH}, \ce{C6H14} (hexane).
\answerlines[3]{\ce{C6H14} $>$ \ce{CH3OH} $>$ \ce{H2O}. Hexane has only
dispersion forces; methanol hydrogen bonds through one O--H; water hydrogen
bonds through two O--H per molecule. Stronger IMFs hold more molecules in
the liquid, lowering vapor pressure.}}

\begin{spiralstrip}{Units 1--2 \textbullet\ spiral}
  \item Geometry and polarity of \ce{SO2}:
        \answerblank[4.5cm]{bent; polar}
  \item PES areas 2:2:6:2 --- element? \answerblank[2.5cm]{Mg}
  \item Rank lattice energy: \ce{NaF}, \ce{MgO}, \ce{KI}:
        \answerblank[4.5cm]{\ce{MgO} $>$ \ce{NaF} $>$ \ce{KI}}
  \item Moles in \SI{11.0}{\gram} \ce{CO2}: \answerblank[2.6cm]{\SI{0.250}{\mole}}
  \item Hybridization of C in \ce{CH4}: \answerblank[2cm]{sp$^3$}
\end{spiralstrip}

\end{document}
